\section{A matched-order memory experiment}
\label{sec:design}

\paragraph{Task and pairing.}
Each episode contains sixteen WRITE observations, a distractor body, and a terminal QUERY (\cref{fig:design}). Eight written coordinates are task-relevant and eight are nuisance. Values are either equiprobable among sixteen categorical levels or independent $\mathrm{Uniform}[-1,1]$ draws. The query selects one required coordinate; neither it nor the body carries the answer. An eight-register oracle is an explicit sufficient-memory control. Paired arms share required values, nuisance values, query, and total updates.

The body uses twelve-symbol motifs over $\{\mathrm A,\mathrm B,\mathrm X\}$, each containing three As, three Bs, and six Xs. The five motifs share bigram counts, unit run lengths, and the prefix \texttt{AX} and suffix \texttt{BX}; the full library is in \cref{app:protocol}. One arm repeats a selected motif $R$ times. The control draws motif identities, fixes its final motif to the same selected motif, and rejects whole-body proper periods. Local repeated substrings remain allowed. At $R=1$, the arms coincide. With unit dwell, the body length is $H=12R$. Matching does not equalize all higher-order correlations or predictability, so the treatment is precisely this motif-ordering intervention.

\paragraph{Acquisition before retention.}
The proposed diagnostic ladder distinguishes a single WRITE followed immediately by its query, all sixteen WRITEs with no body ($R=0$), and short bodies ($R=1,2$). Even $R=0$ requires early values to survive subsequent WRITEs. The stronger design requires every queried coordinate to reach at least $95\%$ categorical accuracy or normalized MSE at most $0.05$ at all three latter horizons in both arms. Normalized MSE divides raw MSE by $1/3$, the continuous population variance. These are proposed criteria: the executed pilots used weaker aggregate gates of $25\%$ accuracy or raw MSE $0.25$ at $R=1,2$, and failed them. Checkpoint selection uses validation task loss, independently of geometric measurements and arm-effect signs.

\paragraph{Behavioral estimands.}
For native error $E_a(R)$, define
\begin{equation}
 \delta_R=E_{\mathrm{per}}(R)-E_{\mathrm{ctl}}(R),\qquad
 \DID=\delta_{16}-\delta_2 .
 \label{eq:contrasts}
\end{equation}
The within-training-support horizons $R=2,16$ compare excess error and its change with retention length. A positive $\DID$ can occur with both arm differences negative; it does not alone establish harm. Independent training seeds are the replication units. Confirmation requires seed-paired uncertainty, a prespecified effect-size threshold and comparison family, and a fixed treatment of competence failures. Ten fresh seeds and a $0.02$ smallest effect are candidate design values, not a completed preregistration. One seed per cell supports no training-population confidence interval. Longer horizons are extrapolation and require separate interpretation.

\paragraph{Localizing a failure.}
Matched episodes are evaluated after WRITE, before QUERY, and after QUERY. Cue-conditioned complete-state clouds separate acquisition, body propagation, and query/readout effects. Train-fit probes use held-out cues, shuffled-label controls, and recorded standardization floors. Offset-compensated decoding tests the alternative in \cref{cor:offset}. State swaps and direction-specific perturbations or rescues test task use against random and nuisance shams; a complete-state swap alone cannot establish mediation by dimension. Restoring an earlier state also changes the operating point, so a rescue requires matched disruption controls. Finite cue differences and robustness to declared noise complement local tangent products. Native precision, coordinate ULP spacing, and fixed-weight FP64 replay are reported separately; replay changes numerical arithmetic and does not retroactively repair stored states.

\paragraph{Decision rule.}
The design permits preserved memory, equal forgetting in both arms, readout failure, and inconclusive outcomes. A learned erasure claim requires prior competence, excess error at a specified horizon, diminished recoverability under the declared measurement model, and evidence that the implicated state change affects behavior beyond shams. A dimension correlation alone does not meet that standard.
